% !TeX root = ../thuthesis-example.tex

\chapter{系统源码与部署说明}

\section{ARM64 精简 Hadoop 镜像构建 Dockerfile 源码}

本项目为解决本机 ARM64 架构兼容性问题，设计并编写了专门的精简版 Hadoop Dockerfile。该文件抛弃了无用组件，只保留 Java 8 与 Hadoop 运行的最基础组件，源码实现如下：
\begin{lstlisting}[style=codexcode]
FROM ubuntu:24.04

# 锁定软件源并拉取基础依赖
RUN apt-get update && apt-get install -y \
    openjdk-8-jdk-headless \
    bash \
    ssh \
    rsync \
    procps \
    curl \
    && rm -rf /var/lib/apt/lists/*

# 设置 JDK 环境变量
ENV JAVA_HOME=/usr/lib/jvm/java-8-openjdk-arm64
ENV PATH=$JAVA_HOME/bin:$PATH

# 导入精简后的 Hadoop 文件包
COPY hadoop /opt/hadoop

ENV HADOOP_HOME=/opt/hadoop
ENV HADOOP_CONF_DIR=/opt/hadoop/etc/hadoop
ENV PATH=/opt/hadoop/bin:/opt/hadoop/sbin:$PATH

# 配置 SSH 密钥实现容器内免密登录
RUN ssh-keygen -t rsa -P '' -f ~/.ssh/id_rsa \
    && cat ~/.ssh/id_rsa.pub >> ~/.ssh/authorized_keys \
    && chmod 0600 ~/.ssh/authorized_keys

WORKDIR /app
\end{lstlisting}

\section{Hadoop 集群核心配置文件 yarn-site.xml 源码}

YARN 容器资源调度由配置驱动，本项目在 \texttt{yarn-site.xml} 中重置了资源上限，防止容器内存过载崩溃。核心调度配置如下：
\begin{lstlisting}[style=codexcode]
<configuration>
    <!-- 指定 ResourceManager 的主机地址 -->
    <property>
        <name>yarn.resourcemanager.hostname</name>
        <value>resourcemanager</value>
    </property>
    
    <!-- 开启辅助服务，为 MapReduce/Spark 提供 Shuffle 存储 -->
    <property>
        <name>yarn.nodemanager.aux-services</name>
        <value>mapreduce_shuffle</value>
    </property>
    
    <!-- 限制单个 NodeManager 分配的最大物理内存 -->
    <property>
        <name>yarn.nodemanager.resource.memory-mb</name>
        <value>2048</value>
    </property>
    
    <!-- 限制单个容器可申请的最大内存 -->
    <property>
        <name>yarn.scheduler.maximum-allocation-mb</name>
        <value>2048</value>
    </property>
    
    <!-- 关闭物理内存与虚拟内存的硬性率上限校验，规避 Docker 环境下的 OOM 误杀 -->
    <property>
        <name>yarn.nodemanager.pmem-check-enabled</name>
        <value>false</value>
    </property>
    <property>
        <name>yarn.nodemanager.vmem-check-enabled</name>
        <value>false</value>
    </property>
</configuration>
\end{lstlisting}

\section{基于 Spark Pipeline 提交与刷新后台任务脚本 submit\_yarn\_client.sh 源码}

本项目专用的 Spark 提交流水线 Shell 脚本。其负责打包代码环境，并将其提报给 YARNResourceManager 执行：
\begin{lstlisting}[style=codexcode]
#!/bin/bash
set -e

# 定位工作根目录
cd "$(dirname "$0")/.."
CONFIG_PATH=$1
RUN_TAG=$2

if [ -z "$CONFIG_PATH" ] || [ -z "$RUN_TAG" ]; then
    echo "Usage: $0 <config_yaml> <run_tag>"
    exit 1
fi

# 打包 PySpark 核心计算任务源码
zip -q -r spark_jobs.zip spark_jobs/ -x "*.pyc" "__pycache__/*"

# 设置环境变量，指向 Docker 内置的 Python 3.11 环境
export PYSPARK_PYTHON=/usr/bin/python3.11
export PYSPARK_DRIVER_PYTHON=/usr/bin/python3.11
export HADOOP_CONF_DIR=/opt/hadoop/etc/hadoop

# 执行 spark-submit 提报
spark-submit \
    --master yarn \
    --deploy-mode client \
    --py-files spark_jobs.zip \
    --name "ecommerce-pipeline-${RUN_TAG}" \
    --driver-memory 1536m \
    --executor-memory 1024m \
    --executor-cores 1 \
    -m spark_jobs.main \
    --config "$CONFIG_PATH" \
    --tag "$RUN_TAG"
\end{lstlisting}

\section{SQLite 数据库初始化与迁移脚本 schema.sql}

项目使用 SQLite 数据库管理异步作业，下面是定义其元数据关系与数据血缘的核心 SQL 脚本：
\begin{lstlisting}[style=codexcode]
-- 任务状态表，记录异步作业元数据
CREATE TABLE IF NOT EXISTS jobs (
    job_id TEXT PRIMARY KEY,
    run_tag TEXT NOT NULL,
    status TEXT NOT NULL,          -- PENDING, RUNNING, SUCCEEDED, FAILED
    error_stage TEXT,              -- 记录出错的 Pipeline 模块
    spark_app_id TEXT,             -- 关联的 YARN Application ID
    created_at TIMESTAMP DEFAULT CURRENT_TIMESTAMP,
    updated_at TIMESTAMP DEFAULT CURRENT_TIMESTAMP
);

-- 数据血缘表，记录输入与输出快照
CREATE TABLE IF NOT EXISTS job_lineage (
    lineage_id INTEGER PRIMARY KEY AUTOINCREMENT,
    job_id TEXT NOT NULL,
    input_snapshot_md5 TEXT,       -- 输入文件的校验码，证明原始一致性
    input_rows INTEGER,
    output_rows INTEGER,
    quality_status TEXT,           -- PASSED, FAILED
    evidence_path TEXT,            -- 指向本地落盘的 run manifest 文件
    FOREIGN KEY(job_id) REFERENCES jobs(job_id)
);
\end{lstlisting}

\section{答辩演示现场验收操作指南}

为了让评审专家在演示现场顺利对系统进行无差错的快速验收，推荐按照如下标准链路进行交互式部署操作：
\begin{enumerate}
    \item **第一步：启动集群基础设施**：
          进入项目根目录，在终端中执行 \texttt{docker compose --profile yarn-lab up -d}。通过 \texttt{docker compose ps} 检查容器状态，确保 NameNode 与 ResourceManager 等 9 个基础节点完全处于正常运行水位。
    \item **第二步：初始化 HDFS 存储**：
          调用初始化命令：\texttt{docker compose --profile yarn-lab exec spark-client /app/scripts/init\_yarn\_lab.sh}。该步骤会在 HDFS 中自动创建必要的分析目录结构，并将本地的 1\% 样本数据物理推送上传至分布式文件系统。
    \item **第三步：提交 Spark 大数据计算任务**：
          在 Spark 客户端容器中执行提交通道脚本：\texttt{docker compose --profile yarn-lab exec spark-client /app/scripts/submit\_yarn\_client.sh /app/configs/yarn-client.yaml demo-run}。
    \item **第四步：启动 Flask 后端与 React 页面服务**：
          宿主机另开终端，运行一键启动脚本 \texttt{./start-yarn-dev.sh}。脚本在清理冲突进程后，启动 API 后端与前端 Vite Web 看板。
    \item **第五步：物理链路追溯验收**：
          在宿主机浏览器中，分别打开 HDFS UI（\texttt{19870}）检查 Parquet 是否分区成功，打开 YARN UI（\texttt{18088}）检查 \texttt{ecommerce-pipeline-demo-run} 的 SUCCEEDED 执行状态，最后打开前端 Web Ops 页面（\texttt{5173/ops}），完成全链路“数理闭环与计算血缘”的工程验收。
\end{enumerate}
